% Scribe section by Anyin Huang
% Slides covered: pages 22--26

\section{Memory Overview}
Taking GPT-3 (175B parameters) as an example, the memory cost of model parameters is about 350 GB (FP16) or 700 GB (FP32). 
The activations, estimated only at transformer layer boundaries for 96 layers, consume about 7488 GB (FP16) or 14976 GB (FP32), though actual internal activations would be higher. 
The optimizer states consume about 12 * 175 GB (FP32) of memory.

\section{Reduce Activation Memory}
In standard training, we store all intermediate activations to compute gradients during the backward pass, leading to an O(N) memory cost for an N-layer network. 
Activation Checkpointing aims to reduce this peak memory by exploiting the fact that activations are not needed until the backward pass. 
Option one is dicard nothing. 
Option two is dicard all and recompute for each layer.
We want to strike a balance. 
The idea is that the activation is not needed again until the backward pass comes, discard some of them and recompute the missing intermediate nodes in small segments. 
For an N-layer neural network, if we place a checkpoint every K layers, the total activation memory cost can be modeled as: Memory Cost = O(N/K) (Checkpoint cost) + O(K) (Re-computation cost). 
By minimizing the cost function, we find the optimal solution at K= $\sqrt{N}$.